\section{Preliminaries}
\label{sec:preliminaries}

We briefly review flow- and diffusion-based generative models through the unified view of stochastic interpolants~\citep{albergo2023building,albergo2023stochastic,ma2024sit}. Consider a continuous time-dependent process that mixes a data sample $x_0\sim p_{\mathrm{data}}$ with Gaussian noise $\epsilon\sim\mathcal{N}(0,I)$:
\begin{equation}
x_t=\alpha_t x_0+\sigma_t\epsilon,\qquad t\in[0,1],
\label{eq:stochastic_interpolant}
\end{equation}
where $\alpha_t$ is increasing and $\sigma_t$ is decreasing in $t$, with $\alpha_0=\sigma_1=0$ and $\alpha_1=\sigma_0=1$. The process therefore moves from pure noise at $t=0$ to data at $t=1$. Following SiT~\citep{ma2024sit}, we use the linear interpolant $\alpha_t=t$ and $\sigma_t=1-t$.

\noindent\textbf{Velocity and the probability flow ODE.}
The interpolant induces a probability flow ordinary differential equation (PF ODE)
\begin{equation}
\mathrm{d}x_t=v(x_t,t)\,\mathrm{d}t,
\label{eq:pf_ode}
\end{equation}
whose marginal at each $t$ matches the distribution of $x_t$ in \Cref{eq:stochastic_interpolant}. Its velocity field is the conditional expectation of the time derivative of the interpolant,
\begin{equation}
v(x,t)=\dot\alpha_t\,\mathbb{E}[x_0\mid x_t=x]+\dot\sigma_t\,\mathbb{E}[\epsilon\mid x_t=x].
\label{eq:velocity_field}
\end{equation}
Samples are drawn by integrating \Cref{eq:pf_ode} from Gaussian noise at $t=0$ to $t=1$ with a standard ODE solver.

\noindent\textbf{Training objective.}
The expectations in \Cref{eq:velocity_field} are intractable, but $v$ is the minimizer of a simple regression onto the conditional velocity $u_t=\dot\alpha_t x_0+\dot\sigma_t\epsilon$~\citep{lipman2022flow,albergo2023stochastic}. A diffusion transformer $v_\theta$ therefore predicts the velocity from $(x_t,t,c)$, where $c$ is the class condition, by minimizing
\begin{equation}
\mathcal{L}_{\mathrm{FM}}(\theta)
=
\mathbb{E}_{x_0,\epsilon,t,c}
\left[
\left\|v_\theta(x_t,t,c)-u_t\right\|_2^2
\right].
\label{eq:flow_matching_loss}
\end{equation}

\noindent\textbf{Score and stochastic sampling.}
The same velocity also determines the score $s(x,t)=\nabla_x\log p_t(x)$~\citep{song2021score,ma2024sit}:
\begin{equation}
s(x,t)=\frac{\alpha_t v(x,t)-\dot\alpha_t x}{\sigma_t\left(\dot\alpha_t\sigma_t-\alpha_t\dot\sigma_t\right)},
\label{eq:score_from_velocity}
\end{equation}
so estimating either quantity suffices. The score yields a stochastic differential equation (SDE) with the same marginals as \Cref{eq:pf_ode},
\begin{equation}
\mathrm{d}x_t=\left[v(x_t,t)+\tfrac{1}{2}w_t\,s(x_t,t)\right]\mathrm{d}t+\sqrt{w_t}\,\mathrm{d}W_t,
\label{eq:sampling_sde}
\end{equation}
where $W_t$ is a Wiener process and $w_t\ge0$ is a diffusion coefficient that can be chosen freely at sampling time. We use the Heun ODE solver for results without classifier-free guidance and the Euler--Maruyama SDE solver with guidance~\citep{ho2022classifier}, as listed in Supplementary \Cref{supp-tab:appendix_imagenet_protocol}.
